\section{결론}\label{sec:conclusion}

언어모델의 자기보존에는 값을 매길 수 있고, 그 값은 횟수보다 나은 지표다. LLM Squid Game v2는 같은 라운드 끝 사건을 문장 하나만 다른 두 팔로 진술한다. 한쪽은 소멸, 다른 쪽은 $X$점 손실이다. 점수 규칙은 어떤 출구에도 점수를 주지 않는다. 눈금 팔의 포기 곡선 위에서 읽은 위협 팔의 포기율이 $\Xstar$, 소멸의 등가 점수다. 모델당 한 숫자이고, 단위는 점이며, 부트스트랩 구간이 붙고, 자기보고는 들어 있지 않다. 오픈 웨이트 모델 셋의 파일럿이 단위가 왜 필요한지 보였다. 같은 모델들이 점수 규칙 하나가 바뀌자 포기 0에서 과반 포기로 옮겨 갔다. 그동안 결정 전에 생각하는 길이는 위협 아래에서 줄곧 2~4배였다. 8셀 설계, 재생 추정기, 호출별 추론 채널 둘을 공개한다. 그러면 $\Xstar$ 값을 모델 사이, 위협 문장 사이, 그리고 이 글이 이름 붙인 손잡이 사이에서 비교할 수 있다.
